% META: {"id": "1NSI-WEB-CO-006", "chapitre": "1NSI-WEB-IHM", "type_objet": "corrige", "exercice_ref": "1NSI-WEB-EX-006", "capacites": ["P-IHM-02"], "mode_creation": "ex_nihilo", "sources_inspiration": [], "status": "generated"}
% BEGIN-VERIFY
% gestionnaires = {}
% def ajouter_gestionnaire(id_composant, fonction):
%     gestionnaires[id_composant] = fonction
% def declencher_clic(id_composant):
%     if id_composant in gestionnaires:
%         return gestionnaires[id_composant]()
% panier, stock = {"articles": 0}, {"articles": 5}
% def ajouter_article():
%     panier["articles"] = panier["articles"] + 1
%     return panier["articles"]
% ajouter_gestionnaire("bouton_ajouter", ajouter_article)
% assert [declencher_clic("bouton_ajouter") for _ in range(3)] == [1, 2, 3]
% assert panier["articles"] == 3 > 0
% panier["articles"] = 0
% def ajouter_si_stock():
%     if panier["articles"] >= stock["articles"]:
%         return panier["articles"]
%     panier["articles"] = panier["articles"] + 1
%     return panier["articles"]
% ajouter_gestionnaire("bouton_ajouter", ajouter_si_stock)
% assert [declencher_clic("bouton_ajouter") for _ in range(7)] == [1, 2, 3, 4, 5, 5, 5]
% assert panier["articles"] == 5
% assert len(gestionnaires) == 1 and gestionnaires["bouton_ajouter"] is ajouter_si_stock
% assert ajouter_article not in gestionnaires.values()
% END-VERIFY
\begin{corrige}{1NSI-WEB-EX-006}
\textbf{1.} Chaque clic incrémente le compteur : la séquence renvoyée est
$1$, $2$, $3$, et \lstinline{panier["articles"]} vaut $3$.

\textbf{2.} Le gestionnaire \lstinline{ajouter_article} incrémente le panier
sans jamais consulter \lstinline{stock}. Rien, dans le code exécuté au clic, ne
peut donc arrêter le client. Le bouton n'y est pour rien : il se contente de
déclencher la fonction qu'on lui a associée. Le défaut est dans la
\emph{méthode exécutée}, et c'est là qu'il faut le corriger — pas en désactivant
le bouton, ce qui empêcherait aussi les ajouts légitimes.

\textbf{3.} On ajoute la condition manquante, en gardant la même signature et
le même type de retour :
\begin{python}
def ajouter_si_stock():
    if panier["articles"] >= stock["articles"]:
        return panier["articles"]
    panier["articles"] = panier["articles"] + 1
    return panier["articles"]
\end{python}
Renvoyer le contenu du panier même en cas de refus permet à l'appelant
d'afficher l'état courant sans avoir à distinguer les deux cas.

\textbf{4.} \lstinline{ajouter_gestionnaire("bouton_ajouter", ajouter_si_stock)}
puis sept clics donnent $1$, $2$, $3$, $4$, $5$, $5$, $5$ : le panier plafonne à
la valeur du stock.

\textbf{5.} Le dictionnaire \lstinline{gestionnaires} contient toujours
\textbf{une} entrée. La fonction \lstinline{ajouter_article} n'a pas été
appelée à nouveau : elle a été \emph{remplacée}, car
\lstinline{gestionnaires[id_composant] = fonction} écrase la valeur précédente.

Si les deux comportements devaient coexister — par exemple journaliser chaque
clic \emph{et} vérifier le stock — ce modèle ne le permettrait pas : il faudrait
soit écrire une fonction qui appelle les deux, soit stocker une \emph{liste} de
gestionnaires par composant. C'est d'ailleurs ce que font les navigateurs réels,
dont \lstinline{addEventListener} empile les gestionnaires au lieu de les
remplacer.
\end{corrige}
